\documentclass[draft,linenumbers]{agujournal2019}
\usepackage{url}
\usepackage{lineno}
\usepackage{soul}

\usepackage{xcolor}
\definecolor{default-linkcolor}{HTML}{A50000}
\definecolor{default-filecolor}{HTML}{A50000}
\definecolor{default-citecolor}{HTML}{4077C0}
\definecolor{default-urlcolor}{HTML}{4077C0}

\usepackage{hyperref}
\hypersetup{
  hidelinks,
  breaklinks=true
}


\providecommand{\tightlist}{\setlength{\itemsep}{0pt}\setlength{\parskip}{0pt}}
\providecommand{\listfigurename}{List of Figures}
\providecommand{\listtablename}{List of Tables}
\newcounter{none}
\setcounter{none}{0}

\makeatletter
\@ifundefined{pandocbounded}{
  \usepackage{graphicx}
  \newcommand{\pandocbounded}[1]{\setbox0=\hbox{#1}\ifdim\wd0>\linewidth\resizebox{\linewidth}{!}{#1}\else#1\fi}
}{}
\makeatother

\usepackage{longtable,booktabs,array}
\usepackage{calc}
\usepackage{etoolbox}
\makeatletter
\patchcmd\longtable{\par}{\if@noskipsec\mbox{}\fi\par}{}{}
\makeatother
\IfFileExists{footnotehyper.sty}{\usepackage{footnotehyper}}{\usepackage{footnote}}
\makesavenoteenv{longtable}

\NewDocumentCommand\citeproctext{}{}
\NewDocumentCommand\citeproc{mm}{\begingroup\def\citeproctext{#2}\cite{#1}\endgroup}
\makeatletter
 \let\@cite@ofmt\@firstofone
 \def\@biblabel#1{}
 \def\@cite#1#2{{#1\if@tempswa , #2\fi}}
\makeatother

\newlength{\cslhangindent}
\setlength{\cslhangindent}{1.5em}
\newlength{\csllabelwidth}
\setlength{\csllabelwidth}{3em}
\newenvironment{CSLReferences}[2]
 {\begin{list}{}{
  \setlength{\itemindent}{0pt}
  \setlength{\leftmargin}{0pt}
  \setlength{\parsep}{0pt}
  \ifodd #1
   \setlength{\leftmargin}{\cslhangindent}
   \setlength{\itemindent}{-\cslhangindent}
  \fi
  \setlength{\itemsep}{#2\baselineskip}
  }
  \item\relax}
 {\end{list}}
\usepackage{calc}
\newcommand{\CSLBlock}[1]{#1\hfill\break}
\newcommand{\CSLLeftMargin}[1]{\parbox[t]{\csllabelwidth}{#1}}
\newcommand{\CSLRightInline}[1]{\parbox[t]{\linewidth - \csllabelwidth}{#1}\break}
\newcommand{\CSLIndent}[1]{\hspace{\cslhangindent}#1}

\usepackage{amsmath}
\usepackage{amssymb}
\usepackage{amsfonts}
\makeatletter
\@ifpackageloaded{subcaption}{}{\usepackage{subcaption}}
\@ifpackageloaded{caption}{}{\usepackage{caption}}
\captionsetup[subfigure]{margin=0.5em}
\AtBeginDocument{%
\renewcommand*\figurename{Figure}
\renewcommand*\tablename{Table}
}
\AtBeginDocument{%
\renewcommand*\listfigurename{List of Figures}
\renewcommand*\listtablename{List of Tables}
}
\newcounter{pandoccrossref@subfigures@footnote@counter}
\newenvironment{pandoccrossrefsubfigures}{%
\setcounter{pandoccrossref@subfigures@footnote@counter}{0}
\begin{figure}\centering%
\gdef\global@pandoccrossref@subfigures@footnotes{}%
\DeclareRobustCommand{\footnote}[1]{\footnotemark%
\stepcounter{pandoccrossref@subfigures@footnote@counter}%
\ifx\global@pandoccrossref@subfigures@footnotes\empty%
\gdef\global@pandoccrossref@subfigures@footnotes{{##1}}%
\else%
\g@addto@macro\global@pandoccrossref@subfigures@footnotes{, {##1}}%
\fi}}%
{\end{figure}%
\addtocounter{footnote}{-\value{pandoccrossref@subfigures@footnote@counter}}
\@for\f:=\global@pandoccrossref@subfigures@footnotes\do{\stepcounter{footnote}\footnotetext{\f}}%
\gdef\global@pandoccrossref@subfigures@footnotes{}}
\@ifpackageloaded{float}{}{\usepackage{float}}
\floatstyle{ruled}
\@ifundefined{c@chapter}{\newfloat{codelisting}{h}{lop}}{\newfloat{codelisting}{h}{lop}[chapter]}
\floatname{codelisting}{Listing}
\newcommand*\listoflistings{\listof{codelisting}{List of Listings}}
\makeatother

\journalname{Geochemistry, Geophysics, Geosystems}

\begin{document}


\title{Along-Strike Heat Flow Continuity Near Subduction Zones Is Mostly
Ambiguous, Not Globally Uniform}

\authors{Buchanan Kerswell\affil{1}, Matthew J. Kohn\affil{2}}

\affiliation{1}{Department of Earth, Oceans and Ecological Sciences,
University of Liverpool, Liverpool L69 3BX, UK}
\affiliation{2}{Department of Geosicences, Boise State University,
Boise, ID 83725}

\correspondingauthor{Buchanan Kerswell}{b.kerswell@liverpool.ac.uk}

\begin{keypoints}
\item Along-strike surface heat flow continuity is geographically
limited; four domains show good continuity, three show clear
discontinuity
\item Kriging and Similarity exhibit complementary failure modes; method
consensus does not guarantee accuracy
\item Nearly 3/4 of all transects had \(\leq\) 25 observations; such
sparse observations are the dominant constraint on geodynamic
interpretation
\end{keypoints}

\begin{abstract}
A common class of subduction zone thermal models assumes slab-mantle
coupling near 70--80 km depth and thin backarc lithospheres, based on
surface heat flow profiles from a small number of well-instrumented
subduction zone segments. These models implicitly assume along-strike
continuity over hundreds of kilometers. We tested this assumption by
applying ordinary Kriging to 60,316 surface heat flow measurements
across 34 subduction zone domains, comparing Krige and Similarity
interpolations point-by-point, and applying cross-correlation function
(CCF) analysis to 201 adjacent trench-perpendicular transect pairs. Four
well-sampled domains (Japan, Cascadia, Sumatra, and Kermadec--Hikurangi)
show sustained continuity consistent with observations and both
interpolation methods. Three other well-sampled domains
(Kurils--Kamchatka, Central America--Panama, and Ryukyus--Sagami) show
abrupt along-strike changes at 100--150 km scales, rendering
single-profile characterizations misleading. Two domains (Caribbean and
Izu--Bonin) exhibit systematic disagreement between methods. Most
transects (168/235; 71\%) have \(\leq\) 25 observations, which is too
sparse for independent assessment. In well-sampled settings, Similarity
mischaracterizes along-strike correlations while Kriging tracks local
observations closely. In sparse settings, only Similarity provides
plausible surface heat flow estimates. Critically, method agreement does
not guarantee accuracy. For Izu--Bonin, 190 local observations
contradict both interpolation methods. For the Caribbean, the two
methods displace a structural boundary by 100--150 km relative to its
inferred location. This analysis shows that the accuracy of thermal
models cannot be determined from the existing surface heat flow record.
Improvement will require strategically targeted surveys and hybrid
Kriging--Similarity approaches guided by the failure-mode diagnostics
defined here.
\end{abstract}

\section*{Plain Language Summary}
When one tectonic plate dives beneath another at a subduction zone, the
two plates lock together and drive flow of warm rock in Earth's mantle.
Scientists measure the heat escaping through Earth's surface to locate
this ``sticking-point'' because its depth controls fault behavior and
earthquake hazards. For decades, a few intensely studied regions were
used to identify a sticking-point depth of 70--80 km, assuming that a
single cross-section could represent the thermal structure for hundreds
of km laterally. We tested this assumption by analyzing more than 60,000
heat flow measurements across 34 subduction zone segments worldwide
using two independent mapping methods. Thermal structure was consistent
in four well-studied regions, supporting the classic model, but thermal
structure changed abruptly over distances of 100--150 km in three other
well-sampled regions, challenging the classic model's assumptions. Two
additional regions gave contradictory results, while nearly three
quarters of all regions we analyzed had too few measurements to evaluate
thermal structure. Importantly, in one region, both methods indicate a
uniform thermal structure while 190 measurements directly contradict
that conclusion; agreement does not prove accuracy. Improving the
subduction zone thermal models requires strategically placed field
surveys combined with mapping strategies that integrate local
measurements and global geological patterns in a principled way.

\clearpage

\section*{Definition of Symbols}\label{sec:symbols}
\addcontentsline{toc}{section}{Definition of Symbols}

{\def\LTcaptype{none} % do not increment counter
\begin{longtable}[]{@{}
  >{\raggedright\arraybackslash}p{(\linewidth - 6\tabcolsep) * \real{0.3481}}
  >{\raggedright\arraybackslash}p{(\linewidth - 6\tabcolsep) * \real{0.2222}}
  >{\raggedright\arraybackslash}p{(\linewidth - 6\tabcolsep) * \real{0.1259}}
  >{\raggedright\arraybackslash}p{(\linewidth - 6\tabcolsep) * \real{0.3037}}@{}}
\toprule\noalign{}
\begin{minipage}[b]{\linewidth}\raggedright
Parameter
\end{minipage} & \begin{minipage}[b]{\linewidth}\raggedright
Symbol
\end{minipage} & \begin{minipage}[b]{\linewidth}\raggedright
Unit
\end{minipage} & \begin{minipage}[b]{\linewidth}\raggedright
Equations
\end{minipage} \\
\midrule\noalign{}
\endhead
\bottomrule\noalign{}
\endlastfoot
\textbf{Spatial Field} & & & \\
Spatial location & \(u\) & km & \ref{eq:variogram}, S1, S3, S4, S5 \\
Observation (surface heat flow) at \(u\) & \(Z(u)\) & mW m\(^{-2}\) &
\ref{eq:variogram}, S1, S5 \\
Mean surface heat flow & \(\mu\) & mW m\(^{-2}\) & S1 \\
Covariance function & \(C(h)\) & mW\(^2\) m\(^{-4}\) & S1, S5 \\
Raw sample count & \(N_\mathrm{obs}\) & - & - \\
\textbf{Variogram Model} & & & \\
Separation distance (lag) between observations & \(h\) & km &
\ref{eq:variogram}, \ref{eq:pairs}, S1, S2, S6 \\
Maximum separation distance & \(\max(h)\) & km & \ref{eq:pairs}, S6 \\
Lag cutoff constant & \(\kappa\) & - & \ref{eq:pairs}, S6 \\
No.~of lag intervals & \(n\) & - & \ref{eq:pairs}, S6, S7 \\
Lag bin width & \(\delta\) & km & S6 \\
Observation pair count at lag \(h\) & \(N(h)\) & - &
\ref{eq:variogram}, \ref{eq:pairs} \\
Experimental variogram & \(\hat{\gamma}(h)\) & mW\(^2\) m\(^{-4}\) &
\ref{eq:variogram} \\
Variogram model & \(\gamma(h)\) & mW\(^2\) m\(^{-4}\) &
\ref{eq:variogram}, S2 \\
Variogram model type & Sph, Exp, Gau & - & S2 \\
Nugget variance & \(\tau^2\) & mW\(^2\) m\(^{-4}\) & S2 \\
Sill variance & \(\sigma^2\) & mW\(^2\) m\(^{-4}\) & S2 \\
Effective range & \(a\) & km & S2 \\
\textbf{Kriging Estimation} & & & \\
Krige estimate at \(u\) & \(\hat{Z}(u)\) & mW m\(^{-2}\) & S3, S4, S5 \\
No.~of observations used for Kriging & \(n_K\) & - & S3, S4, S5 \\
Kriging weight & \(\lambda_i\) & - & S3, S4, S5 \\
Kriging variance & \(\sigma_K^2(u)\) & mW\(^2\) m\(^{-4}\) & S5 \\
Maximum neighborhood observations & \(m\) & - & - \\
Maximum neighborhood size & \(d_\mathrm{max}\) & km & - \\
\textbf{Optimization} & & & \\
Optimization objective & \(J\) & - & S7 \\
Variogram weighted least squares error & \(\mathrm{WLS}_\mathrm{error}\)
& mW\(^4\) m\(^{-8}\) & S7 \\
Standard deviation of experimental variogram & \(\sigma_{\gamma}\) &
mW\(^2\) m\(^{-4}\) & S7 \\
Kriging Root Mean Square Error & \(\mathrm{RMSE}_\mathrm{Krg}\) & mW
m\(^{-2}\) & S7 \\
Standard deviation of cross-validation & \(\sigma_\mathrm{cv}\) & mW
m\(^{-2}\) & S7 \\
\end{longtable}
}

\clearpage

\section{Introduction}\label{sec:introduction}

Surface heat flow reflects the integrated thermal state of the
lithosphere and contributions from hydrothermal circulation, radioactive
decay, fault motion, and mantle convection (Fourier, 1827; Furlong \&
Chapman, 2013; Hutnak et al., 2008; Kelvin, 1863; Stein \& Stein, 1994).
Observations of surface heat flow are therefore central for
understanding lithospheric evolution (Hasterok, 2013; Parsons \&
Sclater, 1977; Pollack \& Chapman, 1977; Rudnick et al., 1998; Stein \&
Stein, 1992), subduction dynamics and seismic hazards (Currie et al.,
2004; Currie \& Hyndman, 2006; Furukawa, 1993; Gao \& Wang, 2014; Kohn
et al., 2018; Wada \& Wang, 2009), groundwater hydrology (Hutnak et al.,
2008; Stein \& Stein, 1994), and exploration of economic resources
including geothermal energy (Batir et al., 2016; Majorowicz \& Grasby,
2010). Decades of community effort have produced tens of thousands of
continental and oceanic measurements compiled into global databases
(Hasterok \& Chapman, 2008; Jennings et al., 2021; Lucazeau, 2019;
Pollack et al., 1993) that underpin multidisciplinary investigations of
lithospheric and crustal processes.

A conventional model of subduction zone thermal structure has emerged
from studies that calibrate numerical thermomechanical models against
individual trench-perpendicular surface heat flow profiles (Currie et
al., 2004; Currie \& Hyndman, 2006; Furukawa, 1993; Hyndman et al.,
2005; Wada \& Wang, 2009), concluding that slab-mantle coupling near
70--80 km depth and a thin backarc lithosphere are common features of
active margins ({Abers et al.}, 2017; Holt \& Condit, 2021; {Syracuse et
al.}, 2010; {van Keken et al.}, 2011, 2018). The model relies on the
implicit but critical assumption that a small number of
well-instrumented transects are spatially representative for hundreds of
kilometers along strike. If proximal transects instead show markedly
different surface heat flow profiles, as we test here, this assumption
breaks down. This raises a fundamental question: do subduction dynamics
genuinely vary over tens to hundreds of kilometers along strike, or do
near-surface perturbations and observational gaps merely obscure a
continuous thermal signature?

Investigating this geodynamic problem requires evaluating transect-scale
surface heat flow patterns across broad regions along strike. The most
recent global compilation from the International Heat Flow Commission
(IHFC 2024) contains 91,182 surface heat flow observations (Global Heat
Flow Data Assessment Group et al., 2024), yet regional coverage near
subduction zones remains sparse and highly uneven at scales of 10--1000
km. High-resolution local surveys, often collocated with seismic
monitoring networks, capture fine-scale variations at individual sites
but do not resolve along-strike variability across subduction zone
segments. The critical methodological challenge is therefore evaluating
sparse, irregularly spaced observations separated by distances ranging
from 0.1--1000 km.

To address both the geodynamic problem and the methodological challenge,
we evaluated two independent interpolation methods that represent
end-member approaches built on fundamentally different principles:
\emph{Kriging} (Krige, 1951) and \emph{Similarity} (Goutorbe et al.,
2011). Their differences are highly consequential for understanding
lithospheric thermal structure, detecting surface heat flow anomalies,
assessing the practical limitations of each approach, and motivating the
development of hybrid methods. The rationale for applying both methods
follows from the \emph{Three Laws of Geography}:

\begin{enumerate}
\def\labelenumi{\arabic{enumi}.}
\tightlist
\item
  Everything is related, but near things are more related than distant
  things (Matheron, 1963)
\item
  Geographic phenomena are inherently heterogeneous (Goodchild, 2004)
\item
  Localities with similar geographic configurations share other
  attributes (Zhu et al., 2018)
\end{enumerate}

Kriging exploits the First Law by estimating surface heat flow in
data-sparse regions using a linear combination of nearby observations,
weighted by their spatial correlations. This reliance on spatial
correlation matches the physics of surface heat flow, which reflects
large-scale lithospheric thermal structure and heat-transfer processes.
For instance, broad regions of low surface heat flow on continents
outline cratons (Nyblade \& Pollack, 1993); anomalously low oceanic
surface heat flow indicates heat extraction by hydrothermally
circulating seawater (Fisher \& Becker, 2000; Hasterok et al., 2011;
Hutnak et al., 2008; Stein \& Stein, 1994); and trench-perpendicular
surface heat flow profiles have been interpreted as evidence for broadly
uniform slab-mantle coupling depths (Furukawa, 1993; Wada \& Wang, 2009)
and upper-plate lithospheric thickness (Currie et al., 2004; Currie \&
Hyndman, 2006; Hyndman et al., 2005).

Similarity, in contrast, does not consider spatial correlations. It
predicts surface heat flow at any target location by comparing that
location's \emph{geographic configuration} (the combination of geologic
and geophysical attributes) against the entire global database of
measured surface heat flow (Goutorbe et al., 2011). This reflects the
Third Law: localities sharing similar bathymetry, lithology, proximity
to orogens, seafloor age, and upper-mantle velocity tend to produce
similar surface heat flow (Chapman \& Pollack, 1975; Davies, 2013; Lee
\& Uyeda, 1965; Lucazeau, 2019; Sclater \& Francheteau, 1970; Shapiro \&
Ritzwoller, 2004). Similarity estimates are therefore geologically
reasoned estimates rather than spatial interpolations.

These considerations motivate a direct test of the assumption underlying
the conventional subduction zone thermal model: that single surface heat
flow profiles are regionally representative of thermal structure, and
that profiles from different subduction systems share common features.
Here we asked three specific questions: a) How spatially continuous is
surface heat flow near active subduction zones? b) How do Krige and
Similarity estimates compare along individual transects and between
adjacent transects? c) What do methodological agreements and
disagreements imply about geodynamic variability and observational gaps?

Our implementation applied Kriging to IHFC 2024 data near 235 predefined
trench-perpendicular transects (organized into 34 transect sets) from
the SubMap tool ({Lallemand et al.}, 2026; Lallemand \& Heuret, 2017),
spanning four major subduction zone regions across the globe. We
compared these Krige estimates point-by-point to global Similarity
estimates from Lucazeau (2019), and used cross-correlation function
(CCF) analysis on adjacent transects to quantify along-strike spatial
continuity. The resulting analysis reveals both region-specific and
method-specific patterns with implications for how surface heat flow
observations constrain subduction zone dynamics.

\section{Methods}\label{sec:methods}

\subsection{Spatial Datasets}\label{sec:datasets}

\subsubsection{IHFC 2024 Surface Heat Flow
Database}\label{sec:ihfc2024-database}

We used the IHFC 2024 database from the Global Heat Flow Data Assessment
Group et al. (2024) for Kriging and CCF analysis. The database contains
91,182 surface heat flow observations from continental and oceanic
settings (Figure~\ref{fig:global-dataset}), together with geographic
coordinates, elevation, unique observation identifiers, and other
metadata. Measurement quality metadata are encoded using the IHFC
uncertainty (U) and measurement (M) classification scheme described by
the Global Heat Flow Data Assessment Group et al. (2024).

We preprocessed the database prior to interpolation. First, we filtered
observations using a simple range criterion, retaining only surface heat
flow values within the interval {[}0, 250{]} mW m\(^{-2}\). This step
excluded non-physical values and extreme observations commonly
associated with localized geothermal systems or measurement artifacts
(Lucazeau, 2019). We applied no additional filtering or weighting based
on IHFC quality classifications.

Second, we identified and resolved duplicate observations located at
identical coordinates. Where duplicate records had different quality
classifications, we retained the observation with the higher combined
uncertainty and measurement quality score (U + M). When duplicate
observations had equal quality scores, we selected one observation
randomly.

After duplicate removal and range filtering, the processed IHFC 2024
dataset contained 60,316 observations used for Kriging and CCF analysis.

\begin{figure}
\centering
\includegraphics[width=0.8\linewidth,height=\textheight,keepaspectratio,alt={Subduction zone segments and global surface heat flow observations. A total of 235 trench-perpendicular transects (top row: bold black lines) and 91,182 raw, unfiltered surface heat flow observations (bottom row) were used for spatial analysis. Convex polygons drawn around 34 sets of transects (top row: shaded areas) defined spatial domains used for sampling and Kriging. Filtered and deduplicated data (N\_\textbackslash mathrm\{obs\} = 60,316 total) show uneven, irregular distributions regionally and globally (Tables S1 and S2). Note the relatively high observation density in N America and around the N Pacific compared to the S Pacific. Plate boundaries (bold white lines) were defined by Coffin et al. (2018). Surface heat flow data are from the Global Heat Flow Data Assessment Group et al. (2024). Transects were defined by Lallemand et al. (2026).}]{../figs/summary/global-dataset.png}
\caption{Subduction zone segments and global surface heat flow
observations. A total of 235 trench-perpendicular transects (top row:
bold black lines) and 91,182 raw, unfiltered surface heat flow
observations (bottom row) were used for spatial analysis. Convex
polygons drawn around 34 sets of transects (top row: shaded areas)
defined spatial domains used for sampling and Kriging. Filtered and
deduplicated data (\(N_\mathrm{obs}\) = 60,316 total) show uneven,
irregular distributions regionally and globally (Tables S1 and S2). Note
the relatively high observation density in N America and around the N
Pacific compared to the S Pacific. Plate boundaries (bold white lines)
were defined by Coffin et al. (2018). Surface heat flow data are from
the Global Heat Flow Data Assessment Group et al. (2024). Transects were
defined by {Lallemand et al.} (2026).}\label{fig:global-dataset}
\end{figure}

\subsubsection{SubMap Transects and Interpolation
Domains}\label{sec:submap-transects}

We used SubMap transects as a standardized geographic framework for
sampling and comparing surface heat flow. SubMap transects are
predefined trench-perpendicular cross-sections through active subduction
zones ({Lallemand et al.}, 2026; Lallemand \& Heuret, 2017). Transects
were organized into 34 geographically contiguous sets spanning four
regional groups: the N Pacific (NPA; 8 sets, 47 transects), S America
(SAM; 10 sets, 66 transects), SE Asia (SEA; 12 sets, 90 transects), and
the SW Pacific (SWP; 4 sets, 32 transects), for a total of 235 transects
(Figure~\ref{fig:global-dataset}; Tables S1 and S2).

For each transect set, we defined an interpolation domain as a convex
polygon enclosing all transects in the set. We sampled IHFC 2024
observations and interpolated surface heat flow grids within each domain
(Figure~\ref{fig:npa-set2-comp}: top row) for CCF analysis. We projected
surface heat flow data onto each transect
(Figure~\ref{fig:npa-set2-comp}: bottom row) by computing the orthogonal
projection distance from each data point to the transect line. A lateral
buffer of 75 km on each side of each transect (150 km total buffer
width) defined the projection window, limiting contributions from points
far from the transect. We sampled a total of 18,635 IHFC 2024
observations within the 34 transect set domains (Table S1), with 9,325
observations located within the projection windows of the 235 transects
(Table S2).

\begin{figure}
\centering
\pandocbounded{\includegraphics[keepaspectratio,alt={Example of surface heat flow interpolation and sampling within the Kurils transect set domain (NPA\_SET2). IHFC 2024 observations (top left) were cropped within a convex polygon (not shown) enclosing the transect set. Krige estimates (top center) were made point-by-point across the same global 0.5\^{}\textbackslash circ \textbackslash times 0.5\^{}\textbackslash circ grid defined by Lucazeau (2019) for Similarity estimates (top right). Both interpolation grids were cropped within the same convex polygon as the IHFC 2024 observations. The 75 km buffers drawn around individual transects (top row: thin black or white rectangles) defined the projection window for the trench-perpendicular profiles used for CCF analysis (bottom row). Plate boundaries (white lines) were defined by Coffin et al. (2018). Surface heat flow data are from the Global Heat Flow Data Assessment Group et al. (2024). Transects were defined by Lallemand et al. (2026).}]{../figs/submap/composition/NPA_SET2-comp.png}}
\caption{Example of surface heat flow interpolation and sampling within
the Kurils transect set domain (NPA\_SET2). IHFC 2024 observations (top
left) were cropped within a convex polygon (not shown) enclosing the
transect set. Krige estimates (top center) were made point-by-point
across the same global 0.5\(^\circ\) \(\times\) 0.5\(^\circ\) grid
defined by Lucazeau (2019) for Similarity estimates (top right). Both
interpolation grids were cropped within the same convex polygon as the
IHFC 2024 observations. The 75 km buffers drawn around individual
transects (top row: thin black or white rectangles) defined the
projection window for the trench-perpendicular profiles used for CCF
analysis (bottom row). Plate boundaries (white lines) were defined by
Coffin et al. (2018). Surface heat flow data are from the Global Heat
Flow Data Assessment Group et al. (2024). Transects were defined by
{Lallemand et al.} (2026).}\label{fig:npa-set2-comp}
\end{figure}

\subsubsection{Map Projection and Interpolation
Grid}\label{sec:map-projection}

We processed spatial datasets using the R package \texttt{sf} (Pebesma,
2018). We imported and stored the IHFC 2024 observations, the Lucazeau
(2019) Similarity interpolation, and the SubMap transects in geographic
coordinates referenced to WGS84 (EPSG:4326). We applied local azimuthal
equidistant projections (ESRI:54032) where required for geometric
buffering operations during construction of convex polygons, then
transformed geometries back to geographic coordinates referenced to
WGS84 (EPSG:4326).

We defined target locations for Kriging using the geometry of the global
surface heat flow grid distributed with the Similarity interpolation of
Lucazeau (2019), enabling direct point-by-point comparison between Krige
estimates and Similarity estimates (e.g., see
Figure~\ref{fig:npa-set2-comp}: top row).

\subsection{Kriging}\label{sec:kriging}

We applied ordinary Kriging to IHFC 2024 observations within each convex
polygon domain using the R package \texttt{gstat} (Gräler et al., 2016;
Pebesma, 2004). Kriging is a three-step process that involves: 1)
computing an experimental variogram \(\hat{\gamma}(h)\) to characterize
spatial correlations, 2) fitting a variogram model \(\gamma(h)\) to the
experimental variogram \(\hat{\gamma}(h)\), and 3) solving a linear
system of equations to predict surface heat flow at target locations
(Cressie, 2015; Krige, 1951; Matheron, 1963). We computed the
experimental variogram as:

\begin{equation}\protect\phantomsection\label{eq:variogram}{
  \hat{\gamma}(h) = \frac{1}{2N(h)} \sum_{i,j \in N(h)} \Big[Z(u_i) - Z(u_j)\Big]^2
}\end{equation}

where \(Z(u_i)\) and \(Z(u_j)\) are observations at locations \(u_i\)
and \(u_j\) separated by lag \(h = \lVert u_i - u_j \rVert\), and
\(N(h)\) is the number of observation pairs separated by that lag
distance. For the irregularly spaced data in this study, we defined
\(N(h)\) using a bin width \(\delta = \max(h)\, /\, \kappa\, n\) such
that:

\begin{equation}\protect\phantomsection\label{eq:pairs}{
  N(h) = \# \left\{(i,j) : h - \frac{\delta}{2} \le h_{ij} < h + \frac{\delta}{2} \right\}
}\end{equation}

where \(\max(h)\) is the maximum separation distance within the Kriging
domain, \(n\) is the number of lag intervals, and \(\kappa\) is the lag
cutoff constant that limits the maximum lag distance \(h\) used to
construct the experimental variogram \(\hat{\gamma}(h)\) (conventionally
\(\kappa\) = 3). We applied Kriging locally by restricting each estimate
to the \(m\) nearest observations that fell within a predefined maximum
neighborhood size \(d_\mathrm{max}\), allowing the effective
interpolation radius to adapt to local data density.

Rather than choosing variogram parameters by eye, a common but
subjective practice, we simultaneously optimized all parameters using
constrained nonlinear optimization (after Li et al., 2018). We tuned the
parameters \(\{n,\, \kappa,\, m,\, d_\mathrm{max}\}\) to minimize a
weighted cost function combining variogram-fitting error and
cross-validation interpolation error. Global optimization used the MLSL
algorithm with Nelder-Mead local refinement. A subsequent fine-tuning
step searched a restricted parameter neighborhood, and we implemented
both stages in the R package \texttt{nloptr}. Full optimization details
and results for all transect sets are presented in the Supplementary
Information (Text S1; Figures S5--S38; Table S9).

\subsection{Cross-Correlation Function
Analysis}\label{sec:cross-correlation-function-analysis}

CCF analysis measured how consistently surface heat flow patterns varied
along strike within each transect set. First, we projected surface heat
flow data from IHFC 2024 observations, Similarity interpolation, and
Krige interpolation onto each transect. To extract continuous
trench-perpendicular trends, we smoothed these projected profiles using
generalized additive models (GAMs, Wood, 2017) via the \texttt{mgcv}
package in R. We modeled smooth profile curves as a function of
along-transect distance using thin plate regression splines, with the
smoothing parameter optimized via Restricted Maximum Likelihood (REML).
We applied the default basis dimension constraint, allowing the
algorithm to dynamically penalize overfitting based on sample size. We
restricted trend fitting to transects containing a minimum of 10
projected data points. Finally, we resampled the smoothed mean profiles
to 500 evenly spaced points along each transect, running from the
subducting plate to approximately 500 km into the upper plate (e.g., see
Figure~\ref{fig:npa-set2-comp}: bottom row).

For each pair of adjacent transects, we quantified the similarity
between surface heat flow profiles using the \texttt{stats::ccf()}
function in R. The analysis allowed small horizontal shifts between
profiles (up to \(\pm\) 10\% of the total transect length) to account
for minor offsets in the positions of surface heat flow anomalies. We
retained the maximum correlation value as a measure of profile
similarity. Correlation coefficients near 1 indicate that adjacent
transects have very similar surface heat flow patterns, whereas values
near 0 or negative indicate weakly related or contrasting patterns.

We performed CCF analysis for all three surface heat flow datasets: IHFC
2024 observations (where data coverage was sufficient), Similarity
interpolation, and Krige interpolation. We computed additional
comparisons between methods (Obs--Sim, Obs--Krg, and Sim--Krg) for
individual transects to evaluate agreement between interpolation
approaches. In total, we analyzed 201 adjacent transect pairs across the
four regional groups (39 in NPA, 56 in SAM, 78 in SEA, and 28 in SWP).

Following the pair-wise calculations, we classified transect sets into
four spatial regimes based on data density and profile correlation
thresholds. Sets with a mean density of fewer than 25 observations per
transect were classified as \textbf{\emph{ambiguous}}. For sets meeting
the 25-observation threshold, classifications were determined by
summarizing cross-method agreement and along-strike continuity:

\begin{itemize}
\tightlist
\item
  \textbf{\emph{Inferred continuous}}: required high method agreement
  within individual transects (mean Obs--Krg CCF \(\geq\) 0.70) and high
  along-strike agreement across all data types (minimum Obs CCF and Krg
  CCF and Sim CCF \(\geq\) 0.50)
\item
  \textbf{\emph{Inferred discontinuous}}: assigned when methods agreed
  within individual transects (mean Obs--Krg CCF \(\geq\) 0.70) but
  adjacent profiles revealed abrupt spatial shifts in either the
  observations (minimum Obs CCF \textless{} 0) or Krige estimates
  (minimum Krg CCF \textless{} 0.30)
\item
  \textbf{\emph{Method disagreement}}: assigned when estimates
  conflicted, either through poor agreement between observations and
  Krige estimates across individual transects (mean Obs--Krg CCF
  \textless{} 0.30) or when observations captured a structural boundary
  (minimum Obs CCF \textless{} 0) that the Similarity method smoothed
  over (maximum Sim CCF \textgreater{} 0.50)
\end{itemize}

Any remaining sets that did not fit these definitive criteria were
designated as \textbf{\emph{ambiguous}}.

\subsection{Interpolation Accuracy}\label{sec:interpolation-accuracy}

We assessed interpolation accuracy by pairing estimated surface heat
flow values from the 0.5\(^\circ\) \(\times\) 0.5\(^\circ\)
interpolation grid directly with the IHFC 2024 observations (after
Lucazeau, 2019). To minimize spatial mismatch, we established point
pairs using a bidirectional nearest-neighbor approach: the nearest
observation to each grid node, and the nearest grid node to each
observation, were cross-referenced. We excluded pairs separated by a
Euclidean distance exceeding a strict 10 km threshold to ensure
localized validity.

The final evaluation dataset mapped each qualified interpolated node to
its nearest valid observation. We summarized estimation errors across
the SubMap transect sets using standard descriptive metrics, including
the Root Mean Square Error (RMSE), mean error, and interquartile range
(IQR).

We applied the same evaluation procedure to both Similarity and Krige
interpolations to ensure consistent comparison between methods. We did
not adopt cross-validation approaches commonly used for Kriging because
the Similarity method uses the full global dataset and is less sensitive
to removal of individual observations (Goutorbe et al., 2011). Direct
comparison with observations therefore provided the most consistent
basis for evaluating both interpolation methods.

\section{Results}\label{sec:results}

\subsection{Along-Strike CCF Analysis Comparing Adjacent Transect
Pairs}\label{sec:along-strike-ccf-analysis-comparing-adjacent-transect-pairs}

CCF analysis across 201 adjacent transect pairs reveals substantial
variation in along-strike surface heat flow continuity among subduction
systems (Figure~\ref{fig:global-ccf-summary}; Table S7). The
classifications shown in Figure~\ref{fig:global-ccf-summary} summarize
the dominant behavior of each transect set rather than requiring perfect
agreement at every transect pair. A classification therefore reflects
the prevailing surface heat flow pattern across a region, while local
departures are examined separately in
Section~\ref{sec:transect-level-ccf-analysis-comparing-observations-kriging-and-similarity}.

Four domains exhibit sustained along-strike continuity over hundreds of
km. Japan (NPA\_SET1), Cascadia (NPA\_SET8), Sumatra (SEA\_SET2), and
Kermadec--Hikurangi (SWP\_SET4) all show moderate-to-high observation
density with strong agreement among observations, Kriging, and
Similarity estimates (Figure~\ref{fig:global-ccf-summary}: green shaded
areas; Table S7). In these regions, neighboring transects consistently
reproduce similar surface heat flow patterns despite spanning large
distances along strike. These domains therefore provide the clearest
evidence that spatially continuous thermal structure can persist across
substantial portions of some subduction segments.

Conversely, three similarly well-sampled domains exhibit abrupt changes
in apparent thermal structure between neighboring transects.
Kurils--Kamchatka (NPA\_SET3), Central America--Panama (SAM\_SET2), and
Ryukyus--Sagami (SEA\_SET9) all contain adjacent transects with weak or
negative correlations that are supported by both observations and Krige
profiles (Figure~\ref{fig:global-ccf-summary}: red shaded areas; Table
S7). In Central America--Panama, strongly correlated transects occur
immediately adjacent to discontinuous pairs, demonstrating that surface
heat flow continuity can change over distances comparable to a single
transect spacing. Manila (SEA\_SET7) displays a similar transition from
continuity to discontinuity, although relatively sparse observations
limit independent validation. Together, these domains indicate that
along-strike thermal structure is not uniformly continuous and may
change abruptly over distances of roughly 100--150 km.

Two domains exhibit mixed continuity patterns that cannot be attributed
confidently to either continuity or discontinuity. The Caribbean
(SAM\_SET4) and Izu--Bonin (SEA\_SET10) both contain alternating
patterns of high and low adjacent-transect similarity, and observations
and interpolation methods support different interpretations of the same
region (Figure~\ref{fig:global-ccf-summary}: yellow shaded areas; Table
S7). These patterns suggest either unresolved thermal heterogeneity at
scales smaller than the interpolation grid or methodological limitations
in reconstructing local surface heat flow structure. Consequently, the
location and significance of apparent thermal regime boundaries remain
uncertain in both domains.

Most transect sets are ambiguous because observation density is
currently too low to reliably evaluate continuity. Transects along large
portions of Alaska, the Caribbean--Andean margin, SE Asia, and the SW
Pacific have few or no observations (Tables S1 and S2), preventing
independent validation of interpolated profiles. In these regions, Krige
estimates are highly sensitive to sparse local data, while Similarity
provides the only spatially complete surface heat flow estimate. As a
result, continuity, discontinuity, and method disagreement cannot be
inferred confidently across 71--75\% of the global subduction zone
system examined here (Table S7).

\begin{figure}
\centering
\pandocbounded{\includegraphics[keepaspectratio,alt={Cross-correlation function (CCF) analysis comparing observation, Krige, and Similarity profiles between adjacent SubMap transect pairs. Maximum CCF values (y-axis) are plotted across 201 adjacent transect pairs arranged by region and set (x-axis), with symbol size scaled by the number of total observations contributing to each pair. Along-strike continuity varies substantially within and across subduction systems. Background shading highlights distinct thermal regimes and method failure modes discussed in the text. Green areas for inferred continuity (NPA\_SET1: Japan; NPA\_SET8: Cascadia; SEA\_SET2: Sumatra; SWP\_SET4: Kermadec--Hikurangi); yellow areas for cross-method disagreement or sub-grid heterogeneity where methods contradict data or give false confidence (SAM\_SET4: Caribbean, SEA\_SET10: Izu--Bonin); red areas for inferred discontinuity along strike (NPA\_SET3: Kurils--Kamchatka; SAM\_SET2: Central America--Panama; SEA\_SET7: Manila; SEA\_SET9: Ryukyus--Sagami); and gray areas for sparse, unconstrained domains where data density is insufficient to infer continuity.}]{../figs/summary/global-ccf-summary.png}}
\caption{Cross-correlation function (CCF) analysis comparing
observation, Krige, and Similarity profiles between adjacent SubMap
transect pairs. Maximum CCF values (y-axis) are plotted across 201
adjacent transect pairs arranged by region and set (x-axis), with symbol
size scaled by the number of total observations contributing to each
pair. Along-strike continuity varies substantially within and across
subduction systems. Background shading highlights distinct thermal
regimes and method failure modes discussed in the text. Green areas for
inferred continuity (NPA\_SET1: Japan; NPA\_SET8: Cascadia; SEA\_SET2:
Sumatra; SWP\_SET4: Kermadec--Hikurangi); yellow areas for cross-method
disagreement or sub-grid heterogeneity where methods contradict data or
give false confidence (SAM\_SET4: Caribbean, SEA\_SET10: Izu--Bonin);
red areas for inferred discontinuity along strike (NPA\_SET3:
Kurils--Kamchatka; SAM\_SET2: Central America--Panama; SEA\_SET7:
Manila; SEA\_SET9: Ryukyus--Sagami); and gray areas for sparse,
unconstrained domains where data density is insufficient to infer
continuity.}\label{fig:global-ccf-summary}
\end{figure}

\begin{figure}
\centering
\pandocbounded{\includegraphics[keepaspectratio,alt={Cross-correlation function (CCF) analysis comparing Obs--Sim (Observed--Similarity), Obs--Krg (Observed--Krige), and Sim--Krg (Similarity--Krige) profile pairs for individual SubMap transects. Maximum CCF values (y-axis) are plotted across individual transects ordered by region and set (x-axis), with symbol size scaled by the number of observations contributing to each transect. High Sim--Krg CCF values identify transects where both interpolation methods recover matching thermal structures, while low or negative values expose systemic method disagreement. Background shading highlights distinct thermal regimes and method failure modes as described in Figure~.}]{../figs/summary/method-ccf-summary.png}}
\caption{Cross-correlation function (CCF) analysis comparing Obs--Sim
(Observed--Similarity), Obs--Krg (Observed--Krige), and Sim--Krg
(Similarity--Krige) profile pairs for individual SubMap transects.
Maximum CCF values (y-axis) are plotted across individual transects
ordered by region and set (x-axis), with symbol size scaled by the
number of observations contributing to each transect. High Sim--Krg CCF
values identify transects where both interpolation methods recover
matching thermal structures, while low or negative values expose
systemic method disagreement. Background shading highlights distinct
thermal regimes and method failure modes as described in
Figure~\ref{fig:global-ccf-summary}.}\label{fig:method-ccf-summary}
\end{figure}

\subsection{Transect-Level CCF Analysis Comparing Observations, Kriging,
and
Similarity}\label{sec:transect-level-ccf-analysis-comparing-observations-kriging-and-similarity}

The classifications in Figure~\ref{fig:global-ccf-summary} describe the
dominant continuity pattern within each transect set, but they do not
imply perfect agreement among observations and interpolation methods at
every transect. Examination of individual transect pairs
(Figure~\ref{fig:method-ccf-summary}) reveals three recurring sources of
uncertainty: 1) disagreement between Similarity and locally constrained
data, 2) Kriging failure despite dense observations, and 3) increasing
interpolation uncertainty in sparsely sampled regions.

The most common disagreement occurs when Similarity produces continuity
estimates inconsistent with both observations and Krige profiles. A
particularly clear example is Ryukyus--Sagami (SEA\_SET9), where
observation density is high (\(N_\mathrm{obs}\) = 83--279) and Obs--Krg
agreement is consistently strong (Obs--Krg CCF = 0.70--1.00;
Figure~\ref{fig:method-ccf-summary}; Table S8). Across five consecutive
transect pairs, observations and Kriging agree on both the sign and
magnitude of adjacent-transect correlations, while Similarity returns a
different interpretation for four of the five pairs
(Figure~\ref{fig:sea67-sea68-sea69-sea70-sea71-sea72-composition}; Table
S7). In several cases, Similarity indicates strong continuity where
observations and Kriging indicate discontinuity, whereas in another pair
it fails to reproduce a nearly perfect correlation identified by both
locally-constrained datasets. Because observations and Kriging converged
across multiple neighboring transects, we attribute the disagreement to
Similarity's failure to resolve the apparent local thermal structure.
Ryukyus--Sagami therefore illustrates a failure mode in which
geologically derived Similarity estimates diverge from observationally
constrained surface heat flow patterns.

The same failure mode appears in several independent subduction systems.
In Central America--Panama (SAM\_SET2), Similarity underestimates
continuity in one transect pair and overestimates it in the immediately
adjacent pair, despite strong support from observations and Kriging for
opposite outcomes. Similar discrepancies occur in Japan (NPA\_SET1),
where Kriging closely tracks observations while Similarity consistently
returns weaker correlations, and in Kurils--Kamchatka (NPA\_SET3), where
Similarity indicates strong continuity even though both observations and
Kriging indicate discontinuity (Figures \ref{fig:global-ccf-summary} and
\ref{fig:method-ccf-summary}; Tables S7 and S8). The recurrence of these
disagreements across multiple regions suggests that Similarity sometimes
fails to reproduce along-strike thermal structure even where local
observations are sufficient to constrain it.

A second failure mode occurs when dense observations do not produce
reliable Krige profiles. In most well-sampled domains, Kriging closely
reproduces observed thermal structure and Obs--Krg CCF values are high.
However, two domains depart from this pattern. In the Caribbean
(SAM\_SET4), the most densely sampled transects in the region (SAM31,
SAM32; \(N_\mathrm{obs}\) = 129, 175) exhibit Obs--Krg CCF values near
zero, indicating that the interpolated Krige profiles bear little
resemblance to the underlying observations
(Figure~\ref{fig:method-ccf-summary}; Table S8). A similar result occurs
in Izu--Bonin (SEA\_SET10), where one densely sampled transect (SEA73;
\(N_\mathrm{obs}\) = 190) also produces near-zero Obs--Krg agreement.
For the adjacent pair SEA73--SEA74, both interpolation methods indicate
positive continuity, yet observations indicate a weak discontinuity
(Table S7). These examples demonstrate that high observation density
alone does not guarantee accurate Krige-based profiles or continuity
estimates.

Interpolation uncertainty increases markedly where observations are
sparse. Point-by-point Sim--Krg prediction differences are centered near
zero in most domains, indicating little systematic bias between methods,
but the spread of those differences varies substantially among regions
(Figure S1; Table S3). The largest discrepancies occur in Andean domains
SAM\_SET9 and SAM\_SET10 and in New Hebrides (SWP\_SET2), where
transects contain few observations and Krige variance is correspondingly
large. Similarity uncertainty remains lower than Krige uncertainty
across most domains (Figure S2; Table S4), whereas in well-sampled
regions Kriging generally produces narrower residual distributions and
stronger agreement with observations (Figures S3 and S4; Tables S5 and
S6). These patterns indicate that uncertainty in surface heat flow
estimates increases systematically as observational constraints weaken.

\begin{figure}
\centering
\pandocbounded{\includegraphics[keepaspectratio,alt={Cross-correlation function (CCF) analysis (top row) and surface heat flow profiles (bottom row) for neighboring Ryukyus--Sagami transects (SEA67--SEA72). Obs and Krg CCF values (left, center) agree in sign and magnitude for all five consecutive transect pairs, even when both are negative, while Similarity CCF (right) returns discordant values in four of the five consecutive pairs (compare first subdiagonals in CCF analysis matrices). Disagreement between observations and Krige estimates versus Similarity reveals Similarity's failure to resolve thermal structure continuity compared to local constraints.}]{../figs/submap/case_study/SEA67-SEA68-SEA69-SEA70-SEA71-SEA72-composition.png}}
\caption{Cross-correlation function (CCF) analysis (top row) and surface
heat flow profiles (bottom row) for neighboring Ryukyus--Sagami
transects (SEA67--SEA72). Obs and Krg CCF values (left, center) agree in
sign and magnitude for all five consecutive transect pairs, even when
both are negative, while Similarity CCF (right) returns discordant
values in four of the five consecutive pairs (compare first subdiagonals
in CCF analysis matrices). Disagreement between observations and Krige
estimates versus Similarity reveals Similarity's failure to resolve
thermal structure continuity compared to local
constraints.}\label{fig:sea67-sea68-sea69-sea70-sea71-sea72-composition}
\end{figure}

\subsection{Optimized Variogram Parameters Across the 34
Domains}\label{sec:optimized-variogram-parameters-across-the-34-domains}

Optimized Krige parameters vary widely among transect sets and show no
evidence for a single characteristic spatial scale of surface heat flow
continuity in subduction zones. All 34 domains converge to an optimal
numerical solution (Figures S5--S38; Table S9). While this consistent
convergence demonstrates the mathematical stability of the optimization
algorithm across a full range of data densities, the resulting
variograms in observation-poor domains (per-transect mean
\(N_\mathrm{obs}\) density \(\leq\) 25) are highly unconstrained and
subject to severe parameter volatility, reflecting localized sample
noise rather than physically meaningful spatial structures.
Nevertheless, the optimization reliably handles the full spectrum of
sampling conditions, from dense networks like Cascadia (up to 1,077
observations per transect) to sparse regions with no immediate transect
observations. Exponential, spherical, and Gaussian variogram models are
each selected across multiple domains; maximum neighborhood sizes
\(d_\mathrm{max}\) range from approximately 100 to 500 km; lag cutoff
constants \(\kappa\) range from 3--12; and the number of lag intervals
\(n\) ranges from 20 to nearly 50 (Figure~\ref{fig:nlopt-summary}; Table
S9). Domains with comparable observation counts frequently converge on
different variogram models and spatial scales, and similarly low-cost
solutions arise from a wide range of parameter configurations. No
systematic relationship exists between individual parameter values and
the total optimization cost.

\begin{figure}
\centering
\includegraphics[width=0.75\linewidth,height=\textheight,keepaspectratio,alt={Summary of optimization iterations, observation counts, and optimized Krige parameters for all transect sets (colored by region). Parameter values (x-axes) span broad ranges and include exponential, spherical, and Gaussian variogram models (Table S9). No systematic relationship exists between individual parameter values and the optimized objective function (y-axis), indicating that characteristic spatial correlation scales vary among subduction zone domains. Optimized Krige parameter units are: maximum neighborhood observations (dimensionless), maximum neighborhood size (km), variogram lag cutoff constant (dimensionless), and number of variogram lag intervals (dimensionless).}]{../figs/summary/nlopt.png}
\caption{Summary of optimization iterations, observation counts, and
optimized Krige parameters for all transect sets (colored by region).
Parameter values (x-axes) span broad ranges and include exponential,
spherical, and Gaussian variogram models (Table S9). No systematic
relationship exists between individual parameter values and the
optimized objective function (y-axis), indicating that characteristic
spatial correlation scales vary among subduction zone domains. Optimized
Krige parameter units are: maximum neighborhood observations
(dimensionless), maximum neighborhood size (km), variogram lag cutoff
constant (dimensionless), and number of variogram lag intervals
(dimensionless).}\label{fig:nlopt-summary}
\end{figure}

\section{Discussion}\label{sec:discussion}

\subsection{Along-Strike Surface Heat Flow Continuity is Ambiguous or
Geographically
Limited}\label{sec:along-strike-heat-flow-continuity-is-ambiguous-or-geographically-limited}

Numerical thermomechanical models of subduction zones have historically
used individual trench-perpendicular surface heat flow profiles to
calibrate their initial thermal structures and boundary conditions. Wada
\& Wang (2009) synthesized models across 17 subduction zones and
concluded that a common slab-mantle coupling depth near 70--80 km
reconciles volcanic arc positions and forearc surface heat flow
globally. Hyndman et al. (2005) and Currie \& Hyndman (2006) identified
broadly uniform backarc surface heat flow of approximately 75 \(\pm\) 15
mW m\(^{-2}\) across ten circum-Pacific systems, attributing this
uniformity to vigorous slab-driven corner flow and a thin, thermally
eroded backarc lithosphere. These conclusions implicitly assume that the
selected surface heat flow profiles all reliably represent thermal
structure and are spatially continuous for hundreds of km along strike.
Our CCF analysis directly evaluates that assumption across 201 adjacent
transect pairs.

Our results support the conventional model of subduction zone thermal
structure only for a minority of the segments where those models were
originally calibrated. Japan (NPA\_SET1), Cascadia (NPA\_SET8), Sumatra
(SEA\_SET2), and Kermadec--Hikurangi (SWP\_SET4) all show sustained
along-strike continuity (Obs CCF = 0.56--1.00), indicating that only
24--40\% of trench-perpendicular profiles underlying the Wada \& Wang
(2009), Hyndman et al. (2005), and Currie \& Hyndman (2006) syntheses
are regionally representative. The tacit assumption that a single
profile captures the broader thermal state of a subduction segment is
therefore geographically limited: elsewhere it either does not hold or
cannot be confirmed. Three well-sampled domains (Kurils--Kamchatka:
NPA\_SET3, Central America--Panama: SAM\_SET2, and Ryukyus--Sagami:
SEA\_SET9) show abrupt along-strike changes at 100--150 km scales, which
render single-profile characterizations misleading. Two additional
domains (Caribbean: SAM\_SET4; Izu--Bonin: SEA\_SET10) show abrupt CCF
transitions whose locations cannot be reliably determined due to method
disagreement, and 71--75\% of the global dataset lacks sufficient
observations for independent assessment (Figures
\ref{fig:global-ccf-summary} and \ref{fig:method-ccf-summary}; Tables S7
and S8). The conventional model of uniform coupling depths with thin
backarc lithospheres is therefore better viewed as a locally validated
description of a small number of well-characterized segments, rather
than as a common feature of active subduction zones worldwide. A nearly
equal number of well-characterized segments show the opposite pattern.

Several deep-seated tectonic mechanisms could explain along-strike
deviations from the conventional model of subduction zone thermal
structure. Numerical geodynamic modeling studies have shown that
variations in slab geometry can concentrate hot subarc mantle at
specific segments via dynamic pressure gradients (Araya Vargas et al.,
2021), and that cessation of backarc spreading can trigger lithospheric
instabilities and along-strike temperature heterogeneity over
million-year timescales (Lee \& Wada, 2021). Kerswell et al. (2021)
showed that backarc lithospheric thickness and serpentinite stability
jointly control the maximum slab-mantle coupling depth, offering another
physical mechanism for the along-strike variability resolved by our CCF
analysis. Peacock \& Wang (2021) proposed a related mineralogical
control involving the pressure-dependent breakdown of talc near 2.0--2.5
GPa (70--80 km depth). Subduction interface friction may also matter:
Kohn et al. (2018) showed that forearc surface heat flow requires shear
heating with friction coefficients of 0.03--0.1, so along-strike
variation in friction {[}surface roughness; Gao \& Wang (2014){]} or
convergence rate could similarly produce the surface heat flow
variability documented here. Where slab geometry, metamorphic
dehydration depth, or interface friction vary along strike, adjacent
profiles can record different thermal regimes within a single subduction
segment.

Yet each of these deeper mechanisms could also, in principle, be
mimicked by near-surface hydrothermal perturbations. Hydrothermal
circulation in young oceanic crust can reduce local heat flow to
10--40\% of lithospheric predictions (Hutnak et al., 2008), and
channelized basement fluid flow can create focused surface heat flow
anomalies over lateral scales of tens of km (Fisher \& Becker, 2000;
Fisher \& Harris, 2010). Without collocated seismic velocity models,
fluid pressure data, or sediment thickness constraints, surface heat
flow cannot discriminate between a slab-mantle coupling depth transition
and a locally perturbed hydrothermal signal. This ambiguity inherently
limits applications of surface heat flow as a geophysical observable in
tectonically active settings.

\subsection{Interpolation Failure Modes Depend on Observation Density
and Proxy
Resolution}\label{sec:interpolation-failure-modes-depend-on-observation-density-and-proxy-resolution}

The three uncertainty patterns identified in our CCF analysis (Sim--Obs
disagreement, Kriging failure in dense settings, and sparsity-driven
uncertainty; see
Section~\ref{sec:transect-level-ccf-analysis-comparing-observations-kriging-and-similarity})
reflect a systematic dependence of method reliability on observation
density. Where per-transect observations are sparse, Krige profiles
overfit individual measurements, leaving Similarity as the only reliable
estimate. Where observations are dense, Kriging closely reproduces local
data while Similarity can mischaracterize along-strike correlations in
specific settings. Understanding when and where each method can fail is
essential for interpreting the surface heat flow record across the range
of data densities represented in the IHFC 2024 dataset (Table S2).

In sparse domains, per-transect observation counts of roughly 25 or
fewer produce Krige profiles sensitive to individual measurements rather
than regional surface heat flow structure. The Aleutian domain
(NPA\_SET5) illustrates this clearly: 4--12 observations per transect
yield highly variable Krg CCF across seven consecutive transect pairs,
while Similarity produces high correlations for six of the seven
(Figure~\ref{fig:global-ccf-summary}; Table S7). In such domains,
Similarity provides the only available geological interpolation.
Although Lucazeau (2019) validated Similarity estimates at the global
scale, systematic validation at the 100--150 km transect-pair scale has
not been previously demonstrated. Whether Similarity accurately captures
along-strike thermal structure within a single subduction segment
therefore remains ambiguous. Where local observations are absent, any
failure of the Similarity method at the transect scale would go
undetected.

In well-sampled domains, the failure-mode pattern reverses and exposes a
specific limitation of Similarity. In Ryukyus--Sagami (SEA\_SET9),
observations and Kriging converge on consistent along-strike CCF values
across five consecutive transect pairs, yet Similarity returns
discordant values for four of the five, including misidentifying
discontinuous segments as continuous and substantially underestimating a
near-perfect correlation inferred by both observations and Krige
estimates
(Figure~\ref{fig:sea67-sea68-sea69-sea70-sea71-sea72-composition}; Table
S7). Central America--Panama (SAM\_SET2) demonstrates that Similarity
failures can occur in both directions within immediately adjacent pairs.
It underestimates strong continuity inferred by observations and Kriging
in one pair, and in another pair returns strong continuity where
observations indicate a structural discontinuity (Figures
\ref{fig:global-ccf-summary} and \ref{fig:method-ccf-summary}; Tables S7
and S8). Similar discrepancies in Japan and Kurils--Kamchatka confirm
that these failures are not confined to a single region.

These systematic failures reflect how Similarity constructs its
estimates. The method identifies global analogs in multivariate proxy
space and applies weighted averaging to their measured values (Goutorbe
et al., 2011). Adjacent transects within a single subduction segment may
share nearly identical proxy configurations while having distinct
along-strike thermal structure. Such differences can arise where proxy
variables (plate age, topography, seismic tomography, and others) fail
to capture important properties or processes, for example fine-scale
slab roughness (interface friction), localized fluid pathways, or
short-wavelength variations in deeper mantle wedge flow. The
0.5\(^\circ\) \(\times\) 0.5\(^\circ\) grid spacing of the Lucazeau
(2019) implementation (approximately 55 km at equatorial latitudes) may
further prevent the proxy field from resolving along-strike thermal
gradients at the transect scale. Prior validation by Shapiro \&
Ritzwoller (2004) demonstrated the viability of proxy-guided
interpolation at regional scales but did not consider the
resolution-dependent limitation identified here.

The Kriging failures in the Caribbean (SAM\_SET4) and Izu--Bonin
(SEA\_SET10) illustrate a complementary limitation. At those transects,
thermal structure varies at spatial scales shorter than any variogram
model's correlation length can resolve. Long variogram ranges
over-smooth the profiles, whereas short ranges are incompatible with
sparse data coverage elsewhere in the domain. No single variogram
configuration reconciles these competing constraints. The convergence of
different Krige parameter configurations across domains therefore
reflects regional differences in the scale of surface heat flow
organization rather than degeneracy in the optimization algorithm (see
Text S1 for full algorithm details). Across the 34 subduction domains
examined here, the absence of a preferred variogram model or
characteristic length scale (Figure~\ref{fig:nlopt-summary}; Table S9)
further suggests that surface heat flow is heterogeneous at all
resolvable scales and effectively unconstrained below a critical
observation-density threshold. A large majority of transects (\(>\)
70\%) currently do not meet this threshold, which is why heat flow
systematics are so ambiguous across much of the globe (Figures
\ref{fig:global-ccf-summary} and \ref{fig:method-ccf-summary}; Tables S7
and S8).

\subsection{Method Agreement Does Not Guarantee
Accuracy}\label{sec:method-agreement-does-not-guarantee-accuracy}

The failure modes described in the previous section share an important
diagnostic signal: method disagreement indicates a problem, but method
agreement does not guarantee accuracy. Two additional patterns in our
analysis make this clear and carry direct implications for how
assessments of subduction zone thermal structure based on surface heat
flow should be used.

The most consequential of these is false consensus between independent
interpolation methods. In Izu--Bonin (SEA\_SET10), Similarity and
Kriging both return positive along-strike continuity (CCF = 0.34--0.93)
for a transect pair (SEA73--SEA74) where 190 local observations indicate
weak discontinuity (CCF = -0.23; Figure~\ref{fig:global-ccf-summary};
Table S7). The Krige profile at the more densely sampled of the two
transects (SEA73) fails to track observed surface heat flow structure
despite the high observation count (Figure~\ref{fig:method-ccf-summary};
Table S8), indicating that local thermal structure varies at scales
shorter than the variogram can resolve. High Sim--Krg CCF agreement
masks this failure rather than flagging it. Neither Similarity's global
proxy resolution nor Kriging's regionally optimized variogram length
scale captures finer-scale heterogeneity, and examining the observations
directly is the only way to detect it. Analyses of this segment that
rely solely on interpolated estimates would find consistent method
agreement supporting a continuity that the measurements directly
contradict.

The Caribbean (SAM\_SET4) documents the distinct complication of
apparent boundary displacement. Observations locate an abrupt thermal
regime boundary at one transect pair (SAM28--SAM29; Obs CCF = -0.73),
while Similarity places it at the immediately adjacent pair
(SAM27--SAM28; Sim CCF = -0.61)---an apparent displacement of roughly
100--150 km (Figure~\ref{fig:global-ccf-summary}). In thermal modeling
contexts where surface heat flow profiles calibrate slab-mantle coupling
depths, a positional error of this magnitude is geodynamically
significant. It misrepresents the lateral extent of the coupled
interface and the along-strike position of key frictional or metamorphic
phase transitions (e.g., Gao \& Wang, 2014; Kohn et al., 2018; Peacock
\& Wang, 2021). The Caribbean's multiple subduction geometries, polarity
reversals, and active backarc basins (Lallemand \& Heuret, 2017)
generate short-wavelength thermal variability that proxy-based methods
cannot resolve at the transect scale, and dense local measurements are
required to accurately locate structural boundaries.

Together, these cases illustrate how method agreement narrows the range
of possible failure modes but does not guarantee accuracy. In
well-constrained domains, agreement between locally-validated Krige
estimates and observationally-consistent Similarity estimates carries
meaningful confidence. Where neither method can be independently checked
against dense local observations, however, agreement may reflect
coincident or correlated failures rather than shared accuracy. The
diagnostic framework developed here, which evaluates per-transect
observation counts alongside CCF analysis, provides a rubric for
distinguishing these situations systematically. Interpolation comparison
without observation-density context is insufficient for interpretation.

\subsection{Sparse Observations Limit Geodynamic
Interpretations}\label{sec:sparse-observations-limit-geodynamic-interpretations}

The transect-level observation counts reported here quantify a practical
threshold for evaluating interpolation reliability. Nearly three
quarters of all transects (168/235; 71\%) lack sufficient observations
for independent validation (\(N_\mathrm{obs}\) \(\leq\) 25), meaning
that the surface heat flow record cannot currently support or refute
geodynamic inferences about along-strike continuity for most of the
global subduction system. This gap reflects a persistent sampling bias
toward accessible continental margins, onshore monitoring networks, and
geothermal exploration targets, leaving subduction zone forearcs and
backarcs systematically undersampled at scales required for geodynamic
interpretation.

Database growth has not yet resolved this issue. The IHFC 2024
compilation represents a substantial addition to prior datasets, adding
roughly 21,000 measurements since the Lucazeau (2019) dataset, yet the
new observations are concentrated in already well-sampled regions. The
along-strike gaps identified here for the Aleutians, most of S America,
and much of SE Asia and the SW Pacific persist, and the \(>\) 1000 km
nearest-neighbor distances documented by Pollack et al. (1993) across
the S Pacific remain relevant for many of the transects examined here.

Where observations do exist, systematic bias in specific tectonic
settings compounds the sparsity problem. Hydrothermal advection accounts
for a substantial fraction of total oceanic heat loss in young oceanic
crust, and conductive measurements in these settings systematically
underestimate the true lithospheric heat flux (Stein \& Stein, 1994).
More than half of the plates entering the trenches analyzed here
(127/235 transects; 54\%) carry oceanic crust younger than the 65
\(\pm\) 10 Ma sealing age beyond which hydrothermal circulation no
longer measurably affects conductive surface heat flow (Stein \& Stein,
1994). Basement fluid discharge along mid-plate seamount chains can
locally reduce regional conductive measurements to a small fraction of
lithospheric predictions (Hutnak et al., 2008). In forearcs and
incoming-plate regions, the available observations are therefore doubly
limited: there are too few measurements to constrain spatial patterns,
and the measurements that do exist are potentially biased by
hydrothermal extraction such that they underestimate the thermal signal
that interpolation methods attempt to reconstruct.

These limitations directly constrain geodynamic inference from the
surface heat flow record. The conventional model of subduction zone
thermal structure proposed by Wada \& Wang (2009) and Currie \& Hyndman
(2006) was built on a small number of well-sampled sites (Japan,
Kamchatka, Cascadia, and Kermadec). Our CCF analysis shows that three of
these sites do exhibit high along-strike continuity, confirming previous
single-profile characterizations as regionally representative. However,
our analysis also identifies one exception: Kurils--Kamchatka
(NPA\_SET3) shows along-strike discontinuity, indicating those profiles
are less representative than previously assumed by Wada \& Wang (2009).
For the broader global system, whether the conventional model extends
beyond the four inferred continuous segments in our CCF analysis
(Figure~\ref{fig:global-ccf-summary}: green shaded areas; Table S7)
cannot be resolved from the existing surface heat flow record
(Figure~\ref{fig:global-ccf-summary}; Table S7). Discontinuities
observed in several segments urge caution.

\subsection{Targeted Surveys and Hybrid
Approaches}\label{sec:targeted-surveys-and-hybrid-approaches}

The data limitations and failure-mode diagnostics discussed above
together define specific, actionable survey priorities rather than a
generic call for more data. Two target categories emerge, each requiring
a different survey strategy and offering a different scientific return.

The most widespread case involves sparse domains where Similarity
provides the only interpolation estimate and no independent validation
is possible. These regions (Figure~\ref{fig:global-ccf-summary}: gray
areas; Table S7) currently have no basis for evaluating whether
Similarity CCF values reflect genuine thermal continuity or simply the
spatial smoothness of the proxy field. Even modest targeted surveys in
these regions would provide the first independent test of Similarity
estimates. The per-transect density thresholds identified here define
the minimum observation density needed to cross from ambiguous to
interpretable (roughly \(N_\mathrm{obs}\) \(>\) 25).

A more demanding strategy is required where fine-scale thermal
heterogeneity exists at scales below global interpolation resolution.
The Caribbean and Izu--Bonin segments demonstrate that short-wavelength
thermal variability can produce abrupt along-strike CCF transitions
within consecutive transect pairs (Figures \ref{fig:global-ccf-summary}
and \ref{fig:method-ccf-summary}; Tables S7 and S8), and that neither
interpolation method reliably locates structural boundaries under these
conditions. Resolving finer-scale structure requires closely spaced
measurements rather than scattered additions. The tectonic complexity of
these systems likely also demands collocated seismic imaging and
geochemical fluid sampling to separate deep thermal signals from shallow
hydrothermal overprinting.

Beyond survey design, the complementary failure modes of Similarity and
Kriging motivate hybrid interpolation strategies. The most direct
combination treats Similarity as a geologically informed prior for the
regional surface heat flow field and applies Kriging to update that
prior with local measurements, conceptually equivalent to Similarity as
a trend model in universal Kriging or as a secondary covariate in a
co-Kriging framework. The failure-mode diagnostics define where each
approach should dominate: Similarity should govern where local
observations are too sparse for Kriging to produce stable estimates,
while Krige estimates should take precedence where per-transect counts
and Obs--Krg CCF indicate strong local constraint. A framework
incorporating the Lucazeau (2019) proxy variables as spatially varying
covariates within a Kriging system would operationalize this without
requiring separate comparison steps.

Kriging and Similarity also perform optimally at different spatial
scales and can be deployed accordingly. The Ryukyus--Sagami analysis
showed that Kriging outperforms Similarity at the transect-pair scale
(Figure~\ref{fig:sea67-sea68-sea69-sea70-sea71-sea72-composition}; Table
S7). A scale-decomposed approach could assign Similarity to characterize
the large-scale background field while Kriging recovers short-scale
departures where local data constrain the variogram. Within the CCF
framework, this decomposition would partition along-strike continuity
patterns between the geological-analog signal in the Similarity field
and the locally constrained spatial structure in Krige residuals, making
failure-mode attribution explicit rather than inferred.

\section{Conclusions}\label{sec:conclusions}

We applied ordinary Kriging to IHFC 2024 data across 235 transects in 34
transect-set domains and quantified along-strike continuity for 201
adjacent transect pairs using cross-correlation function (CCF) analysis.
We compared Krige estimates with the global Similarity interpolation
from Lucazeau (2019) to evaluate where both methods provide independent,
locally constrained assessments of surface heat flow structure. The
combined analysis yields four principal conclusions.

Along-strike surface heat flow continuity is ambiguous or geographically
limited. Four domains (Japan, Cascadia, Sumatra, and
Kermadec--Hikurangi) show sustained continuity consistent with
observations and both interpolation methods, validating common models of
subduction zone thermal structure for those segments (Currie \& Hyndman,
2006; Wada \& Wang, 2009). Three well-sampled domains
(Kurils--Kamchatka, Central America--Panama, and Ryukyus--Sagami) show
abrupt structural changes at 100--150 km along-strike scales independent
of any interpolation method, contradicting assumptions of common
subduction zone thermal models. Two other well-sampled domains (the
Caribbean and Izu--Bonin) exhibit systematic method disagreement with
abrupt along-strike CCF transitions whose locations could not be
confirmed. A large majority of the global dataset lacks sufficient
observations for any independent assessment.

Kriging and Similarity exhibit complementary failure modes. In sparse
domains, Kriging is unreliable and Similarity provides the only
substantive estimate. In well-sampled domains, Kriging tracks
observations closely while Similarity mischaracterizes along-strike
correlations in specific settings. Method consensus does not guarantee
accuracy, but the diagnostic framework developed here provides a basis
for distinguishing reliable from misleading interpretations.

No universal variogram model or characteristic spatial scale describes
subduction zone surface heat flow continuity globally. Kriging
optimization converges for all 34 domains, yet maximum search radii span
approximately 100 to 500 km with no systematic relationship to
optimization cost. Where Krige profiles fail to track local
observations, thermal heterogeneity exists at spatial scales shorter
than any variogram model can resolve. Convergence to different
configurations across domains reflects real regional differences in the
scale of surface heat flow organization rather than algorithmic
degeneracy.

Sparse observations are the dominant limitation on geodynamic
interpretation of subduction across most of the planet. Nearly three
quarters of all transects (168/235; 71\%) have 25 or fewer observations,
precluding Krige-based validation. Additional observations in the IHFC
2024 database relative to prior compilations are concentrated in already
well-sampled regions. In forearc regions, sparse observations may
additionally be biased toward underestimating the true lithospheric heat
flux due to hydrothermal extraction (Hutnak et al., 2008; Stein \&
Stein, 1994). Whether uniform slab-mantle coupling depths and thin
backarc lithospheres are common features of subduction zones remains an
open question. Answering it will require strategically targeted surface
heat flow surveys guided by the failure-mode diagnostics and
observation-density thresholds defined here.

\clearpage

\section*{Acknowledgments}\label{sec:acknowledgements}
\addcontentsline{toc}{section}{Acknowledgments}

We express our sincere gratitude to the communities, institutions, and
research teams who develop and maintain the open-source datasets used in
this study. Specifically, we thank the creators of the SubMap tool
(Géosciences Montpellier); the University of Texas Institute for
Geophysics for the plate boundary compilation; the NOAA National Centers
for Environmental Information for the ETOPO 2022 global relief model;
the Global Heat Flow Data Assessment Group for maintaining the Global
Heat Flow Database; and F. Lucazeau for providing the Similarity
interpolations. This work was supported by the National Science
Foundation grants OIA 1545903 to M. Kohn, S. Penniston-Dorland, and M.
Feineman and EAR 2118114 to M. Kohn.

\section*{Data Availability}\label{sec:data-availability}
\addcontentsline{toc}{section}{Data Availability}

All data, code, and relevant information for reproducing this work are
archived on the OSF (Kerswell, 2026a) and Zenodo (Kerswell, 2026b)
repositories. All code within these repositories is MIT Licensed and
free for use and distribution (see license details). All spatial
datasets used in this study are open-source: transects from the SubMap
tool version 7.1 ({Lallemand et al.}, 2026); present-day plate
boundaries from Coffin et al. (2018); global relief model from NOAA
National Centers for Environmental Information (2022); surface heat flow
observations from the Global Heat Flow Data Assessment Group et al.
(2024); and Similarity interpolations from Lucazeau (2019).

\section*{Conflict of Interest}\label{sec:conflict-of-interest}
\addcontentsline{toc}{section}{Conflict of Interest}

The authors declare there are no conflicts of interest for this
manuscript.

\clearpage

\section*{References}\label{sec:references}
\addcontentsline{toc}{section}{References}

\protect\phantomsection\label{refs}
\begin{CSLReferences}{1}{0}
\bibitem[\citeproctext]{ref-abers2017}
{Abers, G., van Keken, P., \& Hacker, B.} (2017). The cold and
relatively dry nature of mantle forearcs in subduction zones.
\emph{Nature Geoscience}, \emph{10}(5), 333--337.

\bibitem[\citeproctext]{ref-araya2021}
Araya Vargas, J., Sanhueza, J., \& Yáñez, G. (2021). The role of
temperature in the along-margin distribution of volcanism and seismicity
in subduction zones: Insights from 3-d thermomechanical modeling of the
central andean margin. \emph{Tectonics}, \emph{40}(11), e2021TC006879.

\bibitem[\citeproctext]{ref-batir2016}
Batir, J., Blackwell, D., \& Richards, M. (2016). Heat flow and
temperature-depth curves throughout alaska: Finding regions for future
geothermal exploration. \emph{Journal of Geophysics and Engineering},
\emph{13}(3), 366--378.

\bibitem[\citeproctext]{ref-chapman1975}
Chapman, D., \& Pollack, H. (1975). Global heat flow: A new look.
\emph{Earth and Planetary Science Letters}, \emph{28}(1), 23--32.

\bibitem[\citeproctext]{ref-coffin2018}
Coffin, M., Gahagan, L., \& Lawver, L. (2018, July). {Present-day plate
boundary digital data compilation {[}Dataset{]}}. {University of Texas
Institute for Geophysics}. \url{https://doi.org/10.15781/T2ZP3WH84}

\bibitem[\citeproctext]{ref-cressie2015}
Cressie, N. (2015). \emph{Statistics for spatial data}. John Wiley \&
Sons.

\bibitem[\citeproctext]{ref-currie2006}
Currie, C., \& Hyndman, R. (2006). The thermal structure of subduction
zone back arcs. \emph{Journal of Geophysical Research: Solid Earth},
\emph{111}(B8), 1--22.

\bibitem[\citeproctext]{ref-currie2004}
Currie, C., Wang, K., Hyndman, R., \& He, J. (2004). The thermal effects
of steady-state slab-driven mantle flow above a subducting plate: The
cascadia subduction zone and backarc. \emph{Earth and Planetary Science
Letters}, \emph{223}(1-2), 35--48.

\bibitem[\citeproctext]{ref-davies2013}
Davies, J. (2013). Global map of solid earth surface heat flow.
\emph{Geochemistry, Geophysics, Geosystems}, \emph{14}(10), 4608--4622.

\bibitem[\citeproctext]{ref-fisher2000}
Fisher, A., \& Becker, K. (2000). Channelized fluid flow in oceanic
crust reconciles heat-flow and permeability data. \emph{Nature},
\emph{403}(6765), 71--74.

\bibitem[\citeproctext]{ref-fisher2010}
Fisher, A., \& Harris, R. (2010). Using seafloor heat flow as a tracer
to map subseafloor fluid flow in the ocean crust. \emph{Geofluids},
\emph{10}(1-2), 142--160.

\bibitem[\citeproctext]{ref-fourier1827}
Fourier, J. (1827). Mémoire sur les températures du globe terrestre et
des espaces planétaires. \emph{Mémoires de l'Académie Royale Des
Sciences de l'Institut de France}, \emph{7}, 570--604.

\bibitem[\citeproctext]{ref-furlong2013}
Furlong, K., \& Chapman, D. (2013). Heat flow, heat generation, and the
thermal state of the lithosphere. \emph{Annual Review of Earth and
Planetary Sciences}, \emph{41}(1), 385--410.

\bibitem[\citeproctext]{ref-furukawa1993}
Furukawa, Y. (1993). Depth of the decoupling plate interface and thermal
structure under arcs. \emph{Journal of Geophysical Research: Solid
Earth}, \emph{98}(B11), 20005--20013.

\bibitem[\citeproctext]{ref-gao2014}
Gao, X., \& Wang, K. (2014). Strength of stick-slip and creeping
subduction megathrusts from heat flow observations. \emph{Science},
\emph{345}(6200), 1038--1041.

\bibitem[\citeproctext]{ref-ihfc2024}
Global Heat Flow Data Assessment Group et al. (2024, April). {The Global
Heat Flow Database: Release 2024 {[}Dataset{]}} (Version 1). {GFZ Data
Services}. \url{https://doi.org/10.5880/fidgeo.2024.014}

\bibitem[\citeproctext]{ref-goodchild2004}
Goodchild, M. (2004). The validity and usefulness of laws in geographic
information science and geography. \emph{Annals of the Association of
American Geographers}, \emph{94}(2), 300--303.

\bibitem[\citeproctext]{ref-goutorbe2011}
Goutorbe, B., Poort, J., Lucazeau, F., \& Raillard, S. (2011). Global
heat flow trends resolved from multiple geological and geophysical
proxies. \emph{Geophysical Journal International}, \emph{187}(3),
1405--1419.

\bibitem[\citeproctext]{ref-graler2016}
Gräler, B., Pebesma, E., \& Heuvelink, G. (2016). Spatio-temporal
interpolation using gstat. \emph{The R Journal}, \emph{8}, 204--218.
Retrieved from
\url{https://journal.r-project.org/archive/2016/RJ-2016-014/index.html}

\bibitem[\citeproctext]{ref-hasterok2013}
Hasterok, D. (2013). A heat flow based cooling model for tectonic
plates. \emph{Earth and Planetary Science Letters}, \emph{361}, 34--43.

\bibitem[\citeproctext]{ref-hasterok2008}
Hasterok, D., \& Chapman, D. (2008). Global heat flow: A new database
and a new approach. In \emph{AGU fall meeting abstracts} (Vol. 2008, pp.
T21c--1985).

\bibitem[\citeproctext]{ref-hasterok2011}
Hasterok, D., Chapman, D., \& Davis, E. (2011). Oceanic heat flow:
Implications for global heat loss. \emph{Earth and Planetary Science
Letters}, \emph{311}(3-4), 386--395.

\bibitem[\citeproctext]{ref-holt2021}
Holt, A., \& Condit, C. (2021). Slab temperature evolution over the
lifetime of a subduction zone. \emph{Geochemistry, Geophysics,
Geosystems}, \emph{22}(6), e2020GC009476.

\bibitem[\citeproctext]{ref-hutnak2008}
Hutnak, M., Fisher, A., Harris, R., Stein, C., Wang, K., Spinelli, G.,
et al. (2008). Large heat and fluid fluxes driven through mid-plate
outcrops on ocean crust. \emph{Nature Geoscience}, \emph{1}(9),
611--614.

\bibitem[\citeproctext]{ref-hyndman2005}
Hyndman, R., Currie, C., \& Mazzotti, S. (2005). Subduction zone
backarcs, mobile belts, and orogenic heat. \emph{GSA Today},
\emph{15}(2), 4--10.

\bibitem[\citeproctext]{ref-jennings2021}
Jennings, S., Hasterok, D., \& Lucazeau, F. (2021). ThermoGlobe:
Extending the global heat flow database. \emph{Journal TBD}.

\bibitem[\citeproctext]{ref-kelvin1863}
Kelvin, W. (1863). On the secular cooling of the earth.
\emph{Transactions of the Royal Society of Edinburgh}, \emph{23},
157--170.

\bibitem[\citeproctext]{ref-kerswell2026a}
Kerswell, B. (2026a, August).
{buchanankerswell/kerswell\_kohn\_heatflow: Publication Release
{[}Dataset{]}} (Version 1.0.0). OSF.
\url{https://doi.org/10.17605/OSF.IO/CA6ZU}

\bibitem[\citeproctext]{ref-kerswell2026b}
Kerswell, B. (2026b, August).
{buchanankerswell/kerswell\_kohn\_heatflow: Publication Release
{[}Software{]}} (Version 1.0.0). Zenodo.
\url{https://doi.org/10.5281/zenodo.21790947}

\bibitem[\citeproctext]{ref-kerswell2021}
Kerswell, B., Kohn, M., \& Gerya, T. (2021). Backarc lithospheric
thickness and serpentine stability control slab-mantle coupling depths
in subduction zones. \emph{Geochemistry, Geophysics, Geosystems},
\emph{22}(6), e2020GC009304.

\bibitem[\citeproctext]{ref-kohn2018}
Kohn, M., Castro, A., Kerswell, B., Ranero, C., \& Spear, F. (2018).
Shear heating reconciles thermal models with the metamorphic rock record
of subduction. \emph{Proceedings of the National Academy of Sciences},
\emph{115}(46), 11706--11711.

\bibitem[\citeproctext]{ref-krige1951}
Krige, D. (1951). A statistical approach to some basic mine valuation
problems on the witwatersrand. \emph{Journal of the Southern African
Institute of Mining and Metallurgy}, \emph{52}(6), 119--139.

\bibitem[\citeproctext]{ref-lallemand2017}
Lallemand, S., \& Heuret, A. (2017). Subduction zones parameters.
\emph{Reference Module in Earth Systems and Environmental Sciences}.

\bibitem[\citeproctext]{ref-lallemand2026}
{Lallemand, S., Cerpa, N., Peyret, M., Grosbeau, F., Heuret, A., Arcay,
D., \& van Rijsingen, E.} (2026, April). {SubMap: A tool for mapping
subduction zones {[}Software{]}} (Version 7.1). {SubMap}. Retrieved from
\url{https://submap.fr}

\bibitem[\citeproctext]{ref-lee2021}
Lee, C., \& Wada, I. (2021). Volcano clustering promoted by the
cessation of back-arc spreading and ensuing nascent lithospheric drips.
\emph{Geophysical Research Letters}, \emph{48}(9), e2020GL091433.

\bibitem[\citeproctext]{ref-lee1965}
Lee, W., \& Uyeda, S. (1965). Review of heat flow data.
\emph{Terrestrial Heat Flow}, \emph{8}, 87--190.

\bibitem[\citeproctext]{ref-li2018}
Li, Z., Zhang, X., Clarke, K., Liu, G., \& Zhu, R. (2018). An automatic
variogram modeling method with high reliability fitness and estimates.
\emph{Computers \& Geosciences}, \emph{120}, 48--59.

\bibitem[\citeproctext]{ref-lucazeau2019}
Lucazeau, F. (2019). Analysis and mapping of an updated terrestrial heat
flow data set. \emph{Geochemistry, Geophysics, Geosystems},
\emph{20}(8), 4001--4024.

\bibitem[\citeproctext]{ref-majorowicz2010}
Majorowicz, J., \& Grasby, S. (2010). Heat flow, depth--temperature
variations and stored thermal energy for enhanced geothermal systems in
canada. \emph{Journal of Geophysics and Engineering}, \emph{7}(3),
232--241.

\bibitem[\citeproctext]{ref-matheron1963}
Matheron, G. (1963). Principles of geostatistics. \emph{Economic
Geology}, \emph{58}(8), 1246--1266.

\bibitem[\citeproctext]{ref-noaa2022}
NOAA National Centers for Environmental Information. (2022, October).
{ETOPO 2022 15 Arc-Second Global Relief Model {[}Dataset{]}}. {National
Centers for Environmental Information}.
\url{https://doi.org/10.25921/fd45-gt74}

\bibitem[\citeproctext]{ref-nyblade1993}
Nyblade, A., \& Pollack, H. (1993). A global analysis of heat flow from
precambrian terrains: Implications for the thermal structure of archean
and proterozoic lithosphere. \emph{Journal of Geophysical Research:
Solid Earth}, \emph{98}(B7), 12207--12218.

\bibitem[\citeproctext]{ref-parsons1977}
Parsons, B., \& Sclater, J. (1977). An analysis of the variation of
ocean floor bathymetry and heat flow with age. \emph{Journal of
Geophysical Research}, \emph{82}(5), 803--827.

\bibitem[\citeproctext]{ref-peacock2021}
Peacock, S., \& Wang, K. (2021). On the stability of talc in subduction
zones: A possible control on the maximum depth of decoupling between the
subducting plate and mantle wedge. \emph{Geophysical Research Letters},
\emph{48}(17), e2021GL094889.

\bibitem[\citeproctext]{ref-pebesma2004}
Pebesma, E. (2004). Multivariable geostatistics in {S}: The gstat
package. \emph{Computers \& Geosciences}, \emph{30}, 683--691.

\bibitem[\citeproctext]{ref-pebesma2018}
Pebesma, E. (2018). Simple features for r: Standardized support for
spatial vector data. \emph{The R Journal}, \emph{10}(1), 439--446.

\bibitem[\citeproctext]{ref-pollack1977}
Pollack, H., \& Chapman, D. (1977). On the regional variation of heat
flow, geotherms, and lithospheric thickness. \emph{Tectonophysics},
\emph{38}(3-4), 279--296.

\bibitem[\citeproctext]{ref-pollack1993}
Pollack, H., Hurter, S., \& Johnson, J. (1993). Heat flow from the
earth's interior: Analysis of the global data set. \emph{Reviews of
Geophysics}, \emph{31}(3), 267--280.

\bibitem[\citeproctext]{ref-rudnick1998}
Rudnick, R., McDonough, W., \& O'Connell, R. (1998). Thermal structure,
thickness and composition of continental lithosphere. \emph{Chemical
Geology}, \emph{145}(3-4), 395--411.

\bibitem[\citeproctext]{ref-sclater1970}
Sclater, J., \& Francheteau, J. (1970). The implications of terrestrial
heat flow observations on current tectonic and geochemical models of the
crust and upper mantle of the earth. \emph{Geophysical Journal
International}, \emph{20}(5), 509--542.

\bibitem[\citeproctext]{ref-shapiro2004}
Shapiro, N., \& Ritzwoller, M. (2004). Inferring surface heat flux
distributions guided by a global seismic model: Particular application
to antarctica. \emph{Earth and Planetary Science Letters},
\emph{223}(1-2), 213--224.

\bibitem[\citeproctext]{ref-stein1992}
Stein, C., \& Stein, S. (1992). A model for the global variation in
oceanic depth and heat flow with lithospheric age. \emph{Nature},
\emph{359}(6391), 123--129.

\bibitem[\citeproctext]{ref-stein1994}
Stein, C., \& Stein, S. (1994). Constraints on hydrothermal heat flux
through the oceanic lithosphere from global heat flow. \emph{Journal of
Geophysical Research: Solid Earth}, \emph{99}(B2), 3081--3095.

\bibitem[\citeproctext]{ref-syracuse2010}
{Syracuse, E., van Keken, P., \& Abers, G.} (2010). The global range of
subduction zone thermal models. \emph{Physics of the Earth and Planetary
Interiors}, \emph{183}(1-2), 73--90.

\bibitem[\citeproctext]{ref-vankeken2011}
{van Keken, P., Hacker, B., Syracuse, E., \& Abers, G.} (2011).
Subduction factory: 4. Depth-dependent flux of H2O from subducting slabs
worldwide. \emph{Journal of Geophysical Research: Solid Earth},
\emph{116}(B1).

\bibitem[\citeproctext]{ref-vankeken2018}
{van Keken, P., Wada, I., Abers, G., Hacker, B., \& Wang, K.} (2018).
Mafic high-pressure rocks are preferentially exhumed from warm
subduction settings. \emph{Geochemistry, Geophysics, Geosystems},
\emph{19}(9), 2934--2961.

\bibitem[\citeproctext]{ref-wada2009}
Wada, I., \& Wang, K. (2009). Common depth of slab-mantle decoupling:
Reconciling diversity and uniformity of subduction zones.
\emph{Geochemistry, Geophysics, Geosystems}, \emph{10}(10).

\bibitem[\citeproctext]{ref-wood2017}
Wood, S. (2017). \emph{Generalized additive models: An introduction with
{R}}. {Chapman and Hall}.

\bibitem[\citeproctext]{ref-zhu2018}
Zhu, A., Lu, G., Liu, J., Qin, C., \& Zhou, C. (2018). Spatial
prediction based on third law of geography. \emph{Annals of GIS},
\emph{24}(4), 225--240.

\end{CSLReferences}


\end{document}
